\documentclass[11pt,a4paper]{article}
% preamble.tex — estilo común de todos los apuntes Froth.
% Cada apunte hace % preamble.tex — estilo común de todos los apuntes Froth.
% Cada apunte hace % preamble.tex — estilo común de todos los apuntes Froth.
% Cada apunte hace \input{preamble} justo después de \documentclass.
\usepackage[margin=2.4cm]{geometry}
\usepackage{xcolor}
\definecolor{litblue}{RGB}{31,79,121}
\definecolor{litgreen}{RGB}{27,94,32}
\definecolor{litorange}{RGB}{230,81,0}
\usepackage{parskip}
\usepackage{titlesec}
\titleformat{\section}{\Large\bfseries\color{litblue}}{\thesection}{0.6em}{}
\titleformat{\subsection}{\large\bfseries\color{litblue}}{\thesubsection}{0.6em}{}
\usepackage{tcolorbox}
\usepackage{fancyhdr}
\pagestyle{fancy}
\fancyhf{}
\fancyhead[L]{\small\color{litblue}Froth — Apuntes de ML}
\fancyhead[R]{\small\thepage}
\renewcommand{\headrulewidth}{0.4pt}
\usepackage[colorlinks=true,linkcolor=litblue,urlcolor=litblue]{hyperref}

% Cabecera de cada apunte: \cabecera{numero}{titulo}
\newcommand{\cabecera}[2]{%
  \begin{tcolorbox}[colback=litblue,coltext=white,colframe=litblue,arc=2mm]
    {\Large\bfseries Apunte #1 — #2}\\[3pt]
    {\small Proyecto Froth · aprender Machine Learning construyendo}
  \end{tcolorbox}\vspace{0.8em}}

% Cajas temáticas
\newtcolorbox{concepto}[1]{colback=blue!4,colframe=litblue,title={Concepto: #1},arc=1.5mm}
\newtcolorbox{practica}{colback=green!5,colframe=litgreen,title={Pruébalo tú (teclado en mano)},arc=1.5mm}
\newtcolorbox{ojo}{colback=orange!8,colframe=litorange,title={Ojo / error típico},arc=1.5mm}
\newtcolorbox{resumen}{colback=gray!8,colframe=gray!60,title={En una frase},arc=1.5mm}

\newcommand{\codigo}[1]{\texttt{#1}}
 justo después de \documentclass.
\usepackage[margin=2.4cm]{geometry}
\usepackage{xcolor}
\definecolor{litblue}{RGB}{31,79,121}
\definecolor{litgreen}{RGB}{27,94,32}
\definecolor{litorange}{RGB}{230,81,0}
\usepackage{parskip}
\usepackage{titlesec}
\titleformat{\section}{\Large\bfseries\color{litblue}}{\thesection}{0.6em}{}
\titleformat{\subsection}{\large\bfseries\color{litblue}}{\thesubsection}{0.6em}{}
\usepackage{tcolorbox}
\usepackage{fancyhdr}
\pagestyle{fancy}
\fancyhf{}
\fancyhead[L]{\small\color{litblue}Froth — Apuntes de ML}
\fancyhead[R]{\small\thepage}
\renewcommand{\headrulewidth}{0.4pt}
\usepackage[colorlinks=true,linkcolor=litblue,urlcolor=litblue]{hyperref}

% Cabecera de cada apunte: \cabecera{numero}{titulo}
\newcommand{\cabecera}[2]{%
  \begin{tcolorbox}[colback=litblue,coltext=white,colframe=litblue,arc=2mm]
    {\Large\bfseries Apunte #1 — #2}\\[3pt]
    {\small Proyecto Froth · aprender Machine Learning construyendo}
  \end{tcolorbox}\vspace{0.8em}}

% Cajas temáticas
\newtcolorbox{concepto}[1]{colback=blue!4,colframe=litblue,title={Concepto: #1},arc=1.5mm}
\newtcolorbox{practica}{colback=green!5,colframe=litgreen,title={Pruébalo tú (teclado en mano)},arc=1.5mm}
\newtcolorbox{ojo}{colback=orange!8,colframe=litorange,title={Ojo / error típico},arc=1.5mm}
\newtcolorbox{resumen}{colback=gray!8,colframe=gray!60,title={En una frase},arc=1.5mm}

\newcommand{\codigo}[1]{\texttt{#1}}
 justo después de \documentclass.
\usepackage[margin=2.4cm]{geometry}
\usepackage{xcolor}
\definecolor{litblue}{RGB}{31,79,121}
\definecolor{litgreen}{RGB}{27,94,32}
\definecolor{litorange}{RGB}{230,81,0}
\usepackage{parskip}
\usepackage{titlesec}
\titleformat{\section}{\Large\bfseries\color{litblue}}{\thesection}{0.6em}{}
\titleformat{\subsection}{\large\bfseries\color{litblue}}{\thesubsection}{0.6em}{}
\usepackage{tcolorbox}
\usepackage{fancyhdr}
\pagestyle{fancy}
\fancyhf{}
\fancyhead[L]{\small\color{litblue}Froth — Apuntes de ML}
\fancyhead[R]{\small\thepage}
\renewcommand{\headrulewidth}{0.4pt}
\usepackage[colorlinks=true,linkcolor=litblue,urlcolor=litblue]{hyperref}

% Cabecera de cada apunte: \cabecera{numero}{titulo}
\newcommand{\cabecera}[2]{%
  \begin{tcolorbox}[colback=litblue,coltext=white,colframe=litblue,arc=2mm]
    {\Large\bfseries Apunte #1 — #2}\\[3pt]
    {\small Proyecto Froth · aprender Machine Learning construyendo}
  \end{tcolorbox}\vspace{0.8em}}

% Cajas temáticas
\newtcolorbox{concepto}[1]{colback=blue!4,colframe=litblue,title={Concepto: #1},arc=1.5mm}
\newtcolorbox{practica}{colback=green!5,colframe=litgreen,title={Pruébalo tú (teclado en mano)},arc=1.5mm}
\newtcolorbox{ojo}{colback=orange!8,colframe=litorange,title={Ojo / error típico},arc=1.5mm}
\newtcolorbox{resumen}{colback=gray!8,colframe=gray!60,title={En una frase},arc=1.5mm}

\newcommand{\codigo}[1]{\texttt{#1}}

\begin{document}
\cabecera{11}{Tu primera web app con Streamlit}

\section{Qué construimos}

Una \textbf{página web que corre en tu propia PC}: una caja de búsqueda donde escribes una
frase y ves los papers más relevantes en tabla. Es el archivo \codigo{app.py} en la raíz del
proyecto, y por dentro usa el \emph{mismo} \codigo{most\_similar()} de la Fase 2. Cero motor
nuevo: solo una cara bonita encima.

\begin{concepto}{Streamlit}
\textbf{Streamlit} es una librería que convierte un script de Python en una web interactiva
sin que tengas que saber HTML, CSS ni JavaScript. Escribes \codigo{st.title(...)},
\codigo{st.text\_input(...)}, \codigo{st.dataframe(...)} y Streamlit dibuja la página. Cada
vez que el usuario cambia algo (escribe, mueve el slider), Streamlit \textbf{re-ejecuta el
script} de arriba a abajo y actualiza la vista.
\end{concepto}

\section{Cómo arrancarla}

\begin{practica}
Desde la carpeta del proyecto, en PowerShell:
\begin{enumerate}
  \item \codigo{cd "C:\textbackslash Users\textbackslash Bruce Wayne\textbackslash Documents\textbackslash Master Thesis AI"}
  \item \codigo{.venv\textbackslash Scripts\textbackslash python.exe -m streamlit run app.py}
\end{enumerate}
Se abre sola una pestaña en \codigo{http://localhost:8501}. Escribe una consulta, mueve el
slider de resultados y mira la tabla. Para apagarla: \textbf{Ctrl + C} en PowerShell.
\end{practica}

\begin{ojo}
``localhost'' significa ``esta misma máquina''. La web NO está en internet: vive y corre en tu
PC. Por eso no pagas dominio ni hosting. Para que otros la usen: o descargan el repo y la corren
igual, o la subes gratis a Streamlit Community Cloud / Hugging Face Spaces (ver más abajo).
\end{ojo}

\section{Dos trucos del código que vale la pena entender}

\begin{concepto}{caché: no recargar lo pesado en cada clic}
Como Streamlit re-ejecuta el script en cada interacción, cargar los 126 papers y sus vectores
cada vez sería lento. El decorador \codigo{@st.cache\_resource} sobre la función
\codigo{load\_data()} hace que se carguen \textbf{una sola vez} y se reutilicen. El modelo
SPECTER se cachea igual (variable global en \codigo{embed.py}).
\end{concepto}

Además, \codigo{app.py} importa \codigo{from froth import config, embed}: la web es solo un
cliente del motor que ya tienes. Si mañana mejoras \codigo{most\_similar}, la web mejora sola.

\section{Cómo se comparte (tu meta de portafolio)}

\begin{itemize}
  \item \textbf{Descargar y correr:} subes el repo a GitHub; cualquiera hace \codigo{git clone},
        instala \codigo{requirements.txt} y lo corre en su PC. Gratis. (GitHub guarda el código,
        no lo ejecuta por el visitante.)
  \item \textbf{Link público gratis (opcional):} Streamlit Community Cloud o Hugging Face Spaces
        hospedan la app gratis y dan una URL pública. Sin dominio de pago.
\end{itemize}

\section{Lo que viene (visión)}

Esta app es el \textbf{Módulo Tesis}: entrada = un título/frase. Más adelante se añade el
\textbf{Módulo Tarea/Investigación}: adjuntas PDFs de instrucciones + escribes lo que te piden,
y el mismo motor te devuelve los papers relevantes. Y en la Fase 3, además de la lista, aquí
mismo aparecerá el \textbf{bubble map} de subtemas.

\begin{resumen}
Streamlit convierte tu script en una web local (localhost) sin saber HTML. \codigo{app.py}
envuelve \codigo{most\_similar} en una caja de búsqueda; se arranca con \codigo{streamlit run
app.py}, se comparte por GitHub o hosting gratis, y es la base de los dos módulos futuros.
\end{resumen}

\end{document}
